\documentclass[10pt,a4paper]{amsart}
\usepackage{amsmath,amssymb,amsthm,graphicx}
% Retained preamble for the Stacks Project chapter snapshot in this directory.
% The chapter's own preamble.tex is a copy of the shared Stacks Project preamble; the block
% below is sliced MECHANICALLY, line for line, from it. Nothing was added, removed or
% redefined; the wrapper class, the xr-hyper/hyperref wiring and the externaldocument list
% are omitted because the assembler supplies the document class.
%
% lines taken from preamble.tex: 3, 5, 9, 11, 12, 156, 157, 158, ... (51 lines total)
%
% package coverage restored from the chapter's own preamble (2026-10-13): xy, xy-2cell

\usepackage[all]{xy}
\xyoption{2cell}
\UseAllTwocells

\newenvironment{reference}{\comment}{\endcomment}
\newenvironment{slogan}{\comment}{\endcomment}
\newenvironment{history}{\comment}{\endcomment}
\theoremstyle{plain}
\newtheorem{theorem}[subsection]{Theorem}
\newtheorem{proposition}[subsection]{Proposition}
\newtheorem{lemma}[subsection]{Lemma}
\theoremstyle{definition}
\newtheorem{definition}[subsection]{Definition}
\newtheorem{example}[subsection]{Example}
\newtheorem{exercise}[subsection]{Exercise}
\newtheorem{situation}[subsection]{Situation}
\theoremstyle{remark}
\newtheorem{remark}[subsection]{Remark}
\newtheorem{remarks}[subsection]{Remarks}
\numberwithin{equation}{subsection}
\def\lim{\mathop{\mathrm{lim}}\nolimits}
\def\colim{\mathop{\mathrm{colim}}\nolimits}
\def\Spec{\mathop{\mathrm{Spec}}}
\def\Hom{\mathop{\mathrm{Hom}}\nolimits}
\def\Ext{\mathop{\mathrm{Ext}}\nolimits}
\def\SheafHom{\mathop{\mathcal{H}\!\mathit{om}}\nolimits}
\def\SheafExt{\mathop{\mathcal{E}\!\mathit{xt}}\nolimits}
\def\Sch{\mathit{Sch}}
\def\Mor{\mathop{\mathrm{Mor}}\nolimits}
\def\Ob{\mathop{\mathrm{Ob}}\nolimits}
\def\Sh{\mathop{\mathit{Sh}}\nolimits}
\def\NL{\mathop{N\!L}\nolimits}
\def\CH{\mathop{\mathrm{CH}}\nolimits}
\def\proetale{{pro\text{-}\acute{e}tale}}
\def\etale{{\acute{e}tale}}
\def\QCoh{\mathit{QCoh}}
\def\Ker{\mathop{\mathrm{Ker}}}
\def\Im{\mathop{\mathrm{Im}}}
\def\Coker{\mathop{\mathrm{Coker}}}
\def\Coim{\mathop{\mathrm{Coim}}}
\DeclareMathSymbol{\boxtimes}{\mathbin}{AMSa}{"02}
\def\QCohstack{\mathcal{QC}\!\mathit{oh}}
\def\Cohstack{\mathcal{C}\!\mathit{oh}}
\def\Spacesstack{\mathcal{S}\!\mathit{paces}}
\def\Quotfunctor{\mathrm{Quot}}
\def\Hilbfunctor{\mathrm{Hilb}}
\def\Curvesstack{\mathcal{C}\!\mathit{urves}}
\def\Polarizedstack{\mathcal{P}\!\mathit{olarized}}
\def\Complexesstack{\mathcal{C}\!\mathit{omplexes}}
\def\Pic{\mathop{\mathrm{Pic}}\nolimits}
\def\Picardstack{\mathcal{P}\!\mathit{ic}}
\def\Picardfunctor{\mathrm{Pic}}
\def\Deformationcategory{\mathcal{D}\!\mathit{ef}}
% lane-s11 extra preamble for sources whose statements carry \begin{reference} blocks.
% The blocks are purely bibliographic metadata (not mathematical content); the lane
% records the reference in the card's editorial notes, so the body skips them entirely.
\usepackage{comment}
\excludecomment{reference}
\usepackage{cleveref}
\crefname{originalreference}{original reference}{original references}
\Crefname{originalreference}{Original reference}{Original references}
\makeatletter
\newlabel{algebra-lemma-characterize-dvr}{{00PD}{source}{}{sourcepointer-1}{}}
\newlabel{algebra-lemma-characterize-dvr@cref}{{[originalreference][0][]00PD}{[1][0][]source}{}{}{}}
\makeatother
\crefname{originallemma}{original lemma}{original lemmas}
\Crefname{originallemma}{Original lemma}{Original lemmas}
\makeatletter
\newlabel{lemma-morphism-glueing-near-closed-point}{{0GWT}{source}{}{sourcepointer-2}{}}
\newlabel{lemma-morphism-glueing-near-closed-point@cref}{{[originallemma][0][]0GWT}{[1][0][]source}{}{}{}}
\makeatother
\crefname{originallemma}{original lemma}{original lemmas}
\Crefname{originallemma}{Original lemma}{Original lemmas}
\makeatletter
\newlabel{lemma-check-universally-closed-Noetherian}{{05JY}{source}{}{sourcepointer-3}{}}
\newlabel{lemma-check-universally-closed-Noetherian@cref}{{[originallemma][0][]05JY}{[1][0][]source}{}{}{}}
\makeatother
\crefname{originalreference}{original reference}{original references}
\Crefname{originalreference}{Original reference}{Original references}
\makeatletter
\newlabel{morphisms-lemma-finite-type-points-morphism}{{02J3}{source}{}{sourcepointer-4}{}}
\newlabel{morphisms-lemma-finite-type-points-morphism@cref}{{[originalreference][0][]02J3}{[1][0][]source}{}{}{}}
\makeatother
\crefname{originalreference}{original reference}{original references}
\Crefname{originalreference}{Original reference}{Original references}
\makeatletter
\newlabel{morphisms-lemma-chevalley}{{054J}{source}{}{sourcepointer-5}{}}
\newlabel{morphisms-lemma-chevalley@cref}{{[originalreference][0][]054J}{[1][0][]source}{}{}{}}
\makeatother
\crefname{originalreference}{original reference}{original references}
\Crefname{originalreference}{Original reference}{Original references}
\makeatletter
\newlabel{morphisms-lemma-inverse-image-constructible}{{054I}{source}{}{sourcepointer-6}{}}
\newlabel{morphisms-lemma-inverse-image-constructible@cref}{{[originalreference][0][]054I}{[1][0][]source}{}{}{}}
\makeatother
\crefname{originalreference}{original reference}{original references}
\Crefname{originalreference}{Original reference}{Original references}
\makeatletter
\newlabel{morphisms-lemma-ubiquity-Jacobson-schemes}{{02J6}{source}{}{sourcepointer-7}{}}
\newlabel{morphisms-lemma-ubiquity-Jacobson-schemes@cref}{{[originalreference][0][]02J6}{[1][0][]source}{}{}{}}
\makeatother
\crefname{originalreference}{original reference}{original references}
\Crefname{originalreference}{Original reference}{Original references}
\makeatletter
\newlabel{topology-lemma-jacobson-inherited}{{005X}{source}{}{sourcepointer-8}{}}
\newlabel{topology-lemma-jacobson-inherited@cref}{{[originalreference][0][]005X}{[1][0][]source}{}{}{}}
\makeatother
\crefname{originalreference}{original reference}{original references}
\Crefname{originalreference}{Original reference}{Original references}
\makeatletter
\newlabel{morphisms-lemma-closed-point-fibre-locally-finite-type}{{01TF}{source}{}{sourcepointer-9}{}}
\newlabel{morphisms-lemma-closed-point-fibre-locally-finite-type@cref}{{[originalreference][0][]01TF}{[1][0][]source}{}{}{}}
\makeatother
\crefname{originalreference}{original reference}{original references}
\Crefname{originalreference}{Original reference}{Original references}
\makeatletter
\newlabel{schemes-definition-reduced-induced-scheme}{{01J4}{source}{}{sourcepointer-10}{}}
\newlabel{schemes-definition-reduced-induced-scheme@cref}{{[originalreference][0][]01J4}{[1][0][]source}{}{}{}}
\makeatother
\crefname{originalreference}{original reference}{original references}
\Crefname{originalreference}{Original reference}{Original references}
\makeatletter
\newlabel{morphisms-section-normalization-X-in-Y}{{0BAK}{source}{}{sourcepointer-11}{}}
\newlabel{morphisms-section-normalization-X-in-Y@cref}{{[originalreference][0][]0BAK}{[1][0][]source}{}{}{}}
\makeatother
\crefname{originalreference}{original reference}{original references}
\Crefname{originalreference}{Original reference}{Original references}
\makeatletter
\newlabel{morphisms-lemma-finite-type-nagata}{{035A}{source}{}{sourcepointer-12}{}}
\newlabel{morphisms-lemma-finite-type-nagata@cref}{{[originalreference][0][]035A}{[1][0][]source}{}{}{}}
\makeatother
\crefname{originalreference}{original reference}{original references}
\Crefname{originalreference}{Original reference}{Original references}
\makeatletter
\newlabel{morphisms-lemma-nagata-normalization-finite-general}{{0AVK}{source}{}{sourcepointer-13}{}}
\newlabel{morphisms-lemma-nagata-normalization-finite-general@cref}{{[originalreference][0][]0AVK}{[1][0][]source}{}{}{}}
\makeatother
\crefname{originalreference}{original reference}{original references}
\Crefname{originalreference}{Original reference}{Original references}
\makeatletter
\newlabel{morphisms-lemma-rational-map-finite-presentation}{{0BX8}{source}{}{sourcepointer-14}{}}
\newlabel{morphisms-lemma-rational-map-finite-presentation@cref}{{[originalreference][0][]0BX8}{[1][0][]source}{}{}{}}
\makeatother
\title{Testing universal closedness on normal curves punctured at one point (Nagata base)}
\author{The Stacks Project Authors}
\date{stacks-project:limits.tex:local-snapshot-20261012; GNU FDL 1.2 or later}
\begin{document}
\maketitle
\noindent\textit{Editorial note.} Verbatim author omission in the counted span: 'Then we claim that the restriction f' : X' -> Y' of f viewed as a morphism of schemes over S' has property (3) also. We omit the details.' The author flags the paragraph 'we suggest skipping'; it is the mechanical locality of condition (3), while the load-bearing Chevalley/Nagata argument is fully worked. Flagged needs\_human\_spotcheck.\par\medskip
\noindent\textit{Allowed standard prerequisites (editorial).}\begin{itemize}
\item Nagata normalization finiteness. Morphisms, Lemma morphisms-lemma-nagata-normalization-finite-general
\end{itemize}\medskip
\ifdefined\hypertarget\hypertarget{sourcepointer-1}{}\fi
\noindent\textit{Source reference.} Original reference 00PD: sources/stacks-chow/tags.txt, source lines 927--927. Cross-chapter standard prerequisite cited by the author inside the counted span; it is used as a pointer and its own proof is not part of this item.\par\medskip
\ifdefined\hypertarget\hypertarget{sourcepointer-2}{}\fi
\noindent\textit{Source reference.} Original lemma 0GWT: sources/stacks-chow/tags.txt, source lines 21083--21083. Cross-chapter standard prerequisite cited by the author inside the counted span; it is used as a pointer and its own proof is not part of this item.\par\medskip
\ifdefined\hypertarget\hypertarget{sourcepointer-3}{}\fi
\noindent\textit{Source reference.} Original lemma 05JY: sources/stacks-chow/tags.txt, source lines 7182--7182. Cross-chapter standard prerequisite cited by the author inside the counted span; it is used as a pointer and its own proof is not part of this item.\par\medskip
\ifdefined\hypertarget\hypertarget{sourcepointer-4}{}\fi
\noindent\textit{Source reference.} Original reference 02J3: sources/stacks-chow/tags.txt, source lines 3322--3322. Cross-chapter standard prerequisite cited by the author inside the counted span; it is used as a pointer and its own proof is not part of this item.\par\medskip
\ifdefined\hypertarget\hypertarget{sourcepointer-5}{}\fi
\noindent\textit{Source reference.} Original reference 054J: sources/stacks-chow/tags.txt, source lines 6642--6642. Cross-chapter standard prerequisite cited by the author inside the counted span; it is used as a pointer and its own proof is not part of this item.\par\medskip
\ifdefined\hypertarget\hypertarget{sourcepointer-6}{}\fi
\noindent\textit{Source reference.} Original reference 054I: sources/stacks-chow/tags.txt, source lines 6641--6641. Cross-chapter standard prerequisite cited by the author inside the counted span; it is used as a pointer and its own proof is not part of this item.\par\medskip
\ifdefined\hypertarget\hypertarget{sourcepointer-7}{}\fi
\noindent\textit{Source reference.} Original reference 02J6: sources/stacks-chow/tags.txt, source lines 3325--3325. Cross-chapter standard prerequisite cited by the author inside the counted span; it is used as a pointer and its own proof is not part of this item.\par\medskip
\ifdefined\hypertarget\hypertarget{sourcepointer-8}{}\fi
\noindent\textit{Source reference.} Original reference 005X: sources/stacks-chow/tags.txt, source lines 225--225. Cross-chapter standard prerequisite cited by the author inside the counted span; it is used as a pointer and its own proof is not part of this item.\par\medskip
\ifdefined\hypertarget\hypertarget{sourcepointer-9}{}\fi
\noindent\textit{Source reference.} Original reference 01TF: sources/stacks-chow/tags.txt, source lines 2372--2372. Cross-chapter standard prerequisite cited by the author inside the counted span; it is used as a pointer and its own proof is not part of this item.\par\medskip
\ifdefined\hypertarget\hypertarget{sourcepointer-10}{}\fi
\noindent\textit{Source reference.} Original reference 01J4: sources/stacks-chow/tags.txt, source lines 2001--2001. Cross-chapter standard prerequisite cited by the author inside the counted span; it is used as a pointer and its own proof is not part of this item.\par\medskip
\ifdefined\hypertarget\hypertarget{sourcepointer-11}{}\fi
\noindent\textit{Source reference.} Original reference 0BAK: sources/stacks-chow/tags.txt, source lines 14213--14213. Cross-chapter standard prerequisite cited by the author inside the counted span; it is used as a pointer and its own proof is not part of this item.\par\medskip
\ifdefined\hypertarget\hypertarget{sourcepointer-12}{}\fi
\noindent\textit{Source reference.} Original reference 035A: sources/stacks-chow/tags.txt, source lines 4121--4121. Cross-chapter standard prerequisite cited by the author inside the counted span; it is used as a pointer and its own proof is not part of this item.\par\medskip
\ifdefined\hypertarget\hypertarget{sourcepointer-13}{}\fi
\noindent\textit{Source reference.} Original reference 0AVK: sources/stacks-chow/tags.txt, source lines 13688--13688. Cross-chapter standard prerequisite cited by the author inside the counted span; it is used as a pointer and its own proof is not part of this item.\par\medskip
\ifdefined\hypertarget\hypertarget{sourcepointer-14}{}\fi
\noindent\textit{Source reference.} Original reference 0BX8: sources/stacks-chow/tags.txt, source lines 14971--14971. Cross-chapter standard prerequisite cited by the author inside the counted span; it is used as a pointer and its own proof is not part of this item.\par\medskip
\begin{lemma}
\label{lemma-essentially-finite-type-criterion-universally-closed}
Let $S$ be a Nagata scheme (and in particular locally Noetherian).
Let $f : X \to Y$ be a quasi-compact morphism of schemes locally
of finite type over $S$. The following are equivalent
\begin{enumerate}
\item $f$ is universally closed,
\item for every $n$ the morphism
$\mathbf{A}^n \times X \to \mathbf{A}^n \times Y$ is closed,
\item for any commutative diagram
$$
\xymatrix{
U \ar[r] \ar[d] & X \ar[d]^f \\
C \ar[r] \ar@{..>}[ru] & Y
}
$$
of schemes over $S$ such that
\begin{enumerate}
\item $C$ is a normal integral scheme of finite type over $S$,
\item $U = C \setminus \{c\}$ for some closed point $c \in C$,
\item $A = \mathcal{O}_{C, c}$ has dimension $1$\footnote{It follows
that $A$ is a discrete valuation ring, see
Algebra, Lemma \ref{algebra-lemma-characterize-dvr}.
Moreover, $c$ maps to a finite type point $s \in S$ and
$A$ is essentially of finite type over $\mathcal{O}_{S, s}$.}
\end{enumerate}
then in the commutative diagram
$$
\xymatrix{
\Spec(K) \ar[r] \ar[d] & X \ar[d]^f \\
\Spec(A) \ar[r] \ar@{-->}[ru] & Y
}
$$
where $K = \text{Frac}(A)$ some dotted arrow exists\footnote{By
Lemma \ref{lemma-morphism-glueing-near-closed-point} this is
equivalent to asking for the existence of dotted arrow
making the first commutative diagram commute.}
making the diagram commute.
\end{enumerate}
\end{lemma}

\begin{proof}
We have seen the equivalence of (1) and (2) and the
fact that these imply (3) in
Lemma \ref{lemma-check-universally-closed-Noetherian}.
Thus it suffices to prove that (3) implies (2).
Observe that if condition (3) holds for $f : X \to Y$,
then condition (3) holds for
$1 \times f : \mathbf{A}^n \times X \to \mathbf{A}^n \times Y$
(see argument in the proof of
Lemma \ref{lemma-check-universally-closed-Noetherian}).
Hence it suffices to show that (3) implies that $f$ is closed.

\medskip\noindent
Reduction to the case where $Y$ and $S$ are affine; we suggest
skipping this paragraph.
Let $S' \subset S$ be an affine open and let $Y' \subset Y$ be an
affine open mapping into $S'$. Set $X' = f^{-1}(Y')$. Then we
claim that the restriction $f' : X' \to Y'$ of $f$ viewed as a morphism
of schemes over $S'$ has property (3) also. We omit the details.
Now if we can prove that $f'$ is closed for all choices of
$S'$ and $Y'$, then it follows that $f$ is closed. This reduces
us to the case discussed in the next paragraph.

\medskip\noindent
Assume $S$ and $Y$ affine. Let $Z \subset X$ be a closed subset.
We may and do view $Z$ as a reduced closed subscheme of $X$.
We have to show that $E = f(Z)$ is closed. Pick $y \in Y$ a closed point
contained in the closure of $f(Z)$. It suffices to show $y \in E$.
We assume $y \not \in E$ to get a contradiction.
The image $s \in S$ of $y$ is a finite type point of $S$, see
Morphisms, Lemma \ref{morphisms-lemma-finite-type-points-morphism}.
Recall that $E$ is constructible
(Morphisms, Lemma \ref{morphisms-lemma-chevalley}).
Consider the intersection $\Spec(\mathcal{O}_{Y, y}) \cap E$.
This is a constructible subset of the spectrum
(Morphisms, Lemma \ref{morphisms-lemma-inverse-image-constructible})
which doesn't contain the closed point.
Since the punctured spectrum
$\Spec(\mathcal{O}_{Y, y}) \setminus \{y\}$ is Jacobson
(Morphisms, Lemma \ref{morphisms-lemma-ubiquity-Jacobson-schemes}), we find a
closed point $t \in \Spec(\mathcal{O}_{Y, y}) \setminus \{y\}$
with $t \in E$ (see Topology, Lemma \ref{topology-lemma-jacobson-inherited}).
In other words, $t \in E$ is a point of $Y$ which
has an immediate specialization $t \leadsto y$.
As $t \in E$ the scheme theoretic fibre $Z_t$ is nonempty.
Choose a closed point $x \in Z_t$. In particular we have
$[\kappa(x) : \kappa(t)] < \infty$ by the Hilbert Nullstellensatz
(Morphisms, Lemma
\ref{morphisms-lemma-closed-point-fibre-locally-finite-type}).

\medskip\noindent
Denote $T = \overline{\{t\}} \subset Y$ the integral closed subscheme
whose underlying topological space is as indicated
(Schemes, Definition \ref{schemes-definition-reduced-induced-scheme}).
Then $t \in T$ is the generic point.
Denote $C \to T$ the normalization of $T$ in $\kappa(x)$, see
Morphisms, Section \ref{morphisms-section-normalization-X-in-Y}
(more precisely, $C \to T$ is the normalization of $T$ in $x$
where we view $x = \Spec(\kappa(x)) \to T$ as a scheme over $T$).
Since $S$ is a Nagata scheme, so is $T$
(Morphisms, Lemma \ref{morphisms-lemma-finite-type-nagata}).
Hence we see that $C \to T$ is finite
(Morphisms, Lemma \ref{morphisms-lemma-nagata-normalization-finite-general}).
As $t$ is in the image we see that $C \to T$ is surjective (because
the image is closed and $T$ is the closure of $t$ in $Y$).
Choose a point $c \in C$ mapping to $y \in T$.
Since $y$ is a closed point of $T$ we see that $c$ is
a closed point of $C$. Since $\dim(\mathcal{O}_{T, y}) = 1$
we see that $\dim(\mathcal{O}_{C, c}) = 1$ (the dimension is
at least $1$ as $c$ is not the generic point of $C$ and
at most $1$ as $C \to T$ is finite).
As the function field of $C$ is $\kappa(x)$ and as $x$
is a point of $X$, we have a $Y$-rational map from $C$ to $X$
(see for example Morphisms, Lemma
\ref{morphisms-lemma-rational-map-finite-presentation}).
Let $C \supset U \to X$ be a representative (in particular $U$ is nonempty).
We may assume $c \not \in U$ (replace $U$ by $U \setminus \{c\}$).
Since $c$ is a closed point of codimension $1$ in the integral scheme $C$
we have $C = U \amalg \{c\} \amalg \Sigma$
for some proper closed subset $\Sigma \subset C$.
After replacing $C$ by $C \setminus \Sigma$
we have constructed a commutative diagram as in part (3).
By the 2nd footnote in the statement of the lemma,
the existence of the dotted arrow produces an extension
of the rational map to all of $C$ and we get the contradiction
because the image of $c$ will be a point of $Z$ mapping to $y$.
\end{proof}

\end{document}
